\reviewsection{When a Label Opens a Door---and When It Starts Closing One}

In my elementary technology teaching placement, I repeatedly saw students whose competence was easy to miss because of what surrounded the task. A learner might understand the coding concept but struggle with a long verbal direction, an unfamiliar login path, reading load, noise, peer control of the device, or uncertainty about what finished work should look like. When I modeled the exact tool path, used a visible first/next/last sequence, broke the task into smaller steps, or gave additional processing time, participation changed. The academic idea had not suddenly entered the student's mind. The environment had stopped hiding what was already there.

That experience is what made \citet{cole2023autism} meaningful to me. Their clinical and educational language can give families and professionals a shared vocabulary for social communication, sensory processing, routines, focused interests, and support needs. The CDC and broader school-centered overview similarly validate the practical value of diagnostic literacy: identification can open eligibility, services, and coordinated support \citep{cdcASD,simon2016asd}. At the same time, all three sources acknowledge heterogeneity. A diagnostic label cannot tell me what a particular learner understands, fears, enjoys, prefers, or needs in a specific moment. It cannot tell me whether the barrier is sensory, linguistic, technological, social, emotional, or simply a confusing set of directions.

The wider source set also keeps this claim from becoming too comfortable. \citet{nadesan2013genetics} asks how genetic knowledge becomes entangled with profit, risk, policy, and judgments about which lives are desirable. \citet{orsiniDavidson2013critical} likewise challenge epidemic language and the false choice between treating autism only as a disorder or only as a difference. These sources do not ask educators to discard medical knowledge. They ask us to notice what medical and scientific narratives do once they enter schools, families, funding systems, and public imagination.

This is why labels can become dangerous when adults let them replace curiosity. They influence what people expect, which behaviors are interpreted as meaningful, and which students are assumed capable before they have a chance to show it. Diagnosis is one source of evidence, but it is not the whole story. The challenge is not to reject diagnostic knowledge; it is to prevent that knowledge from becoming the only knowledge adults treat as credible.

\reviewsection{Neurodiversity Became Real to Me Through Both Identity and Teaching}

\citet{ortega2013cerebralizing} connected with something I have felt for years: labels are not only clinical. They can become part of identity, community, resistance, and self-understanding. As someone who received an ADHD diagnosis later in life, I know the relief of finding language that explains patterns previously interpreted as laziness, inconsistency, or failure. A name can reduce shame. It can help a person reinterpret the past and ask for support without believing they are morally defective.

But Ortega also warns that an affirming brain-based identity can become reductive if it suggests that a person is only a brain or that everyone with the same label shares one political, cultural, or personal experience. That tension appears in classrooms too. A teacher can move away from deficit language and still create a new stereotype: the autistic child who is always visual, the student with ADHD who is always energetic, or the neurodivergent learner who is automatically gifted in one narrow way.

My students keep reminding me that difference does not arrive in neat categories. In one activity, a student needed a familiar laptop nearby while also building with blocks. The device could have been interpreted as distraction or refusal. Instead, I treated it as part of the student's regulation and allowed both forms of participation to coexist. The student remained connected to the activity.

\citet{conn2016play} helped me see why that moment mattered beyond accommodation. Conn challenges the familiar myth that autistic children do not truly play and questions adult-directed approaches that measure whether a child performs the outward signs of normative play. The reading directly complicates any intervention that values visible compliance more than internal engagement, enjoyment, or relationship. It also reports that autistic children may sustain more interaction when partners are less directive and allow more time. That does not mean adults withdraw support; it means support must leave room for the learner's authorship.

The Autism Speaks reading list is useful as a pathway to family-facing and public-facing resources, but even a resource list carries assumptions about whose books, voices, and purposes count \citep{autismSpeaksReads}. It therefore needs to be paired with autistic-authored, critical, and neurodiversity-affirming work. For me, neurodiversity is most useful when it expands whose knowledge counts and protects variation within the spectrum. It should make us more curious, not more certain.

\reviewsection{What the Media Helped Me Recognize About Representation}

Watching \textit{Autism: Insight from Inside} reminded me of how strongly public understanding depends on stories \citep{infobase2016inside}. I have experienced the power of being seen through the wrong story: the unreliable student, the overly emotional person, or the one who should simply try harder. I also know the relief of encountering a story that makes room for complexity rather than forcing a person into a deficit narrative.

The film validates several claims in the readings. Its participants describe labels as both access points and sources of stigma, and they insist that people with the same diagnosis can have very different sensory experiences, communication profiles, interests, and support needs. At the same time, parts of the film complicate my own argument. Its strong emphasis on work skills, independence, productivity, and ``good outcomes'' can unintentionally imply that acceptance must be earned through economic usefulness. Building strengths is important, but a person's dignity cannot depend on becoming employable, independent, or exceptional.

The broader media menu makes this tension even clearer. \textit{Autism: The Musical} shows children creating, performing, relating, and surprising adults, which strongly validates the claim that environments can reveal competence \citep{regan2007musical}. Yet its older dialogue also includes pressure to ``blend in'' and family movement through denial, fixing, and acceptance. The course warning that the film is nearly two decades old matters because representation can affirm possibility while still carrying normalization assumptions from its time.

\textit{Best Kept Secret} extends the question beyond childhood by following autistic students in Newark as they approach the end of public-school services \citep{buck2013best}. It supports my claim that disability is socially interpreted, but it also forces a harder question: what does school-based inclusion mean if students encounter a service cliff afterward? Its attention to race, economic inequality, and adult transition strengthens the intersectional argument that access is never distributed evenly.

The course also offers \textit{Talk to Me} and \textit{Unspoken: The Dynamic Life of an Autistic Teenager} \citep{talkToMe2007,unspoken2018}. Because the licensed versions were not locally transcribed, I do not assign them detailed claims that I cannot verify. Still, their placement in this module matters: communication and adolescence are not side topics but central sites where adults decide whose expression is recognizable and whose life is considered knowable. The American Autism Association podcast pathway further broadens the media field through parent, teaching, sensory, faith-based, Black autism, and first-person podcast perspectives \citep{aaa2022podcasts}. Its range validates the need for many voices, while its strong concentration of parent- and professional-centered shows also reminds me to ask whether autistic speakers are being heard directly.

Media can widen public understanding when it centers multiple voices and makes context visible. It can narrow understanding when one inspirational, tragic, productive, or heroic story becomes the measuring stick for everyone else.
